% !TeX program = pdflatex
% ALIUS Bulletin native reconstruction source.
% Compile from the ALIUS-bulletin project root. This file does not include a pre-existing article PDF.
\providecommand{\ALIUSRootPrefix}{}

\ifdefined\ALIUSIssueBuild
\else
\documentclass[a4paper,11pt]{article}
\usepackage[margin=22mm]{geometry}
\usepackage[T1]{fontenc}
\usepackage[utf8]{inputenc}
\usepackage[hidelinks]{hyperref}
\usepackage[protrusion=true,expansion=false]{microtype}
\usepackage{parskip}
\fi
\title{Consciousness : from molecular to social scales}
\author{An interview with Jean-Pierre Changeux by Guillaume Dumas}
\date{ALIUS Bulletin 7}

\ifdefined\ALIUSIssueBuild
\else
\begin{document}
\fi
\maketitle
\section*{Metadata}
\textbf{Piece type:} interview\par
\textbf{DOI:} 10.5281/zenodo.13907411\par
\textbf{Keywords:} consciousness, neuromodulation, global neuronal workspace, AI, art, psychedelics\par
\section*{Abstract}
In this interview, Guillaume Dumas reaches out to Jean-Pierre Changeux to offer a wide-ranging view on consciousness, spanning over its fundamental molecular mechanisms up to its social dimension. They further discuss the exploration of altered states of consciousness and its relation to artistic practice. They finally tackle the opportunities of some of the prominent fields of future research such as artificial intelligence and highlight the importance of multidisciplinary research in advancing the study of consciousness.\par
\section*{Extracted Text}
\begingroup
\small
\noindent Consciousness : from molecular to social scales. An interview with Jean-Pierre Changeux By Guillaume Dumas Citation: Changeux, J-P. \& Dumas, G. (2024). Consciousness: from molecular to social scales. An interview with Jean-Pierre Changeux. ALIUS Bulletin, XX, XX-XX Jean-Pierre Changeux changeux@noos.fr Institut Pasteur, Paris, France Guillaume Dumas guillaume.dumas@umontreal.ca Universite de Montreal, Montreal, Canada Abstract In this interview, Guillaume Dumas reaches out to Jean-Pierre Changeux to offer a wide-ranging view on consciousness, spanning over its fundamental molecular mechanisms up to its social dimension. They further discuss the exploration of altered states of consciousness and its relation to artistic practice. They finally tackle the opportunities of some of the prominent fields of future research such as artificial intelligence and highlight the importance of multidisciplinary research in advancing the study of consciousness. Keywords: consciousness, neuromodulation, global neuronal workspace, AI, art, psychedelics Your work has significantly influenced our understanding of the neural underpinnings of consciousness. How do you view the evolution of consciousness studies since your pioneering forays into the field? First of all, in the past decades, an authentic scientific approach to conscious processing was initiated by Francis Crick with his "spotlight of attention" hypothesis about visual awareness. His provocative position paradoxically gave legitimacy to the reintegration of the word consciousness into contemporary scientific language but also offered the opportunity to develop a theoretical and experimental understanding of consciousness. ALIUS Bulletin n7 (2024) aliusresearch.org/bulletin Changeux - Consciousness: from molecular to social scales 2 From the start, I would like to strongly underline the importance of theory, which is not often the case, in those days, in biological and physiological sciences, at variance with physics, where theory is dominant. However, in our discipline, theory without experimental tests does not make sense. Fortunately, in the field of consciousness research, new experimental tools became available, in the past decades, in brain imaging, genetics, molecular biology, chemistry, and AI... which - together with the joint use of experimental psychology - made possible a fair evaluation of theoretical hypotheses about consciousness right away. An "integrative" strategy based on the embodiment of cognitive functions, establishing a causal relationship between neuro-molecular processes, modeling, and conscious processing, became accessible. I dislike the word "neural correlates" of consciousness but prefer "neural mechanisms" or, better "neuro-molecular mechanisms" of consciousness. The future is there, together with a particular emphasis on the strong relationships established by the human brain with its multiple environments biological, physical, social, and cultural. Concerning consciousness itself, up to now, it has been investigated almost exclusively with the individual brain and at its simplest levels. A very promising future is opened, with a similar "adapted" strategy, for the social brain and the higher levels of consciousness including theory of mind, rationality, language... A fascinating avenue! Throughout your career, what discoveries or insights have most altered your own understanding of consciousness? I would prefer not to use "altered" but "progress" in my understanding of consciousness. In my last book ("Le Beau et la Splendeur du Vrai") I have described my paths from allostery, "Neuronal man", up to the Global Neuronal Workspace. Early on, in Neuronal Man, I had a section (Ch 5.9) on the "substance of the spirit" where I introduced the idea that a "surveillance system" composed of very divergent neurons and their reentries, such that a "spider web" (of connections), may form and function as a "whole" and as a consequence consciousness would "emerge." ALIUS Bulletin n7 (2024) aliusresearch.org/bulletin Changeux - Consciousness: from molecular to social scales 3 The teaching at the College de France gave me the opportunity to further reflect on the fascinating history of the concepts relative to consciousness in the course of the past centuries, starting from Lamarck's "sentiment interieur," Hughling Jackson's hierarchical organization, William James's stream of consciousness, Llinas's electrophysiology of conscious states and last but not least Baars's book "a cognitive theory of consciousness" (1988). In this book, Baars introduces the Global Workspace theory (GW), yet with little, if any, solid neurobiological foundations and no reference to the molecular level. Beforehand, together with Stanislas Dehaene, I had elaborated neuro-molecular models of low-level cognitive functions such as the Wisconsin Card Sorting Task (Dehaene \& Changeux 1991) and after reading Baars's, I decided to extend the same approach to the neurobiological and molecular mechanisms involved in the GW. We proposed, among other features, the contribution of a brain scale network of cortical neurons with long-range axons, spontaneous activity of the network, dopamine reward... as basic components of a Global "Neuronal" Workspace (GNW) (Dehaene, Kerszberg, Changeux 1998; Dehaene \& Changeux, 2011). This was the beginning of an amazing story that is still quite alive! In your opinion, how do psychoactive substances modulate neuronal networks, and what implications does this have for our understanding of consciousness? Drugs modulate neuronal networks at the level of circuits of neurons releasing a well-defined neurotransmitter. They are intrinsic components of the circuits engaged in any defined brain functions. For example, dopamine is a neurotransmitter well known to play a fundamental role in reward, and cocaine/amphetamines block dopamine re-uptake by its presynaptic transporter. The issue is then: would some of these neuro-molecular circuits contribute to conscious processing? where and how? The original version of the Global Neuronal Workspace of 1998 was designed to account for the "Stroop task" and the reward was an essential part of the global network (as recently emphasized by Volzhenin et al., 2022). ALIUS Bulletin n7 (2024) aliusresearch.org/bulletin Changeux - Consciousness: from molecular to social scales 4 Tasks subsequently used to test the GNW, like the "masking tasks", did not need explicit reward. Interestingly, the modeling of a masking task (Dehaene et al., 2003) led to the proposal of a contribution to two critical components of brain circuitry: the NMDA vs AMPA glutamate receptors. These would intervene differentially in top-down vs bottom-up "resonant" circuits engaged in conscious processing. Further modeling and experimental work confirmed that. Sooner or later, the fair understanding of conscious processing shall obligatorily include the molecular level! Could you elaborate on the neurobiological mechanisms that you believe are most crucial in altering states of consciousness, particularly in relation to your research on nicotinic receptors and allosteric modulation? Most often, the current studies of the states of consciousness in well-doing subjects rely upon behavior, clinics, or computational modeling. "Disorders" of consciousness like coma, vegetative state, or minimally conscious states... reveal the importance of critical elements of brain anatomy in the access to consciousness. What is usually reported as "altered" states of consciousness like, for instance, the hallucinations produced by psychotropic drugs, are especially interesting for me since they selectively target the molecular level on top of which the whole brain is built. There is a fundamental bottom-up top-down relationship between the molecular and the conscious level, which is frequently underestimated in the field. A wide variety of pharmacological agents, natural or synthetic, are known to act differentially on brain molecules such as neurotransmitter receptors but also enzymes, transporters. Many of them bind directly to the "active" site and have the same or a more powerful effect as the neurotransmitter - they are referred to as agonists. Others block the physiological ligand at the same site in a mutually exclusive manner - they are referred to as competitive antagonists. Nicotine at high doses is used by Amazonian populations to "shape dreams" and by Australian Aboriginal people to cause trance-like states and "excite their courage in warfare". It is a potent agonist of the nicotinic acetylcholine receptor. ALIUS Bulletin n7 (2024) aliusresearch.org/bulletin Changeux - Consciousness: from molecular to social scales 5 There is a fundamental bottom-up top-down relationship between the molecular and the conscious level, which is frequently underestimated in the field. The beloved Martin Fortier (from the Ecole Normale Superieure) distinguished two main categories of hallucinogenic drugs: first, the serotonergic drugs like mushroom psilocybin, whose action resembles that of LSD which behaves as an agonist of the serotonin 5-HT2A receptors, and second, the muscarinic drugs, which are competitive antagonists of the G-protein linked muscarinic receptor for acetylcholine. In relation to our research on nicotinic receptors, it was discovered that some drugs do not act at the level of the receptor binding site nor at the level of the biologically active site (channel, G-protein site...) but target different allosteric modulatory sites distributed at different levels on the receptor molecule. One example is the benzodiazepines with potent anxiolytic properties which act as positive allosteric modulators of the GABAA receptor. Another example is related to the effects of hashish: tetrahydrocannabinol and cannabinoid ligands act as positive allosteric modulators of glycine receptors. About 80 allosteric modulators are used in the clinic (Changeux \& Christopoulos, 2016)! A new allosteric pharmacology is born to investigate the disorders of consciousness. Considering your extensive research in neuropharmacology, what ethical considerations should guide us in exploring altered states of consciousness through substances? This is a necessary but challenging question for me as the former Chair of the French National Bioethics Commission. My first reaction is that because of the diversity of drugs and their mode of action, each case should be examined and discussed by itself. The second reaction is that the subject's informed consent should be received with a fair explanation of the experiment's aim, which is most often difficult to really explain to the patient. The third is that there exists in many countries special committees for the protection of persons that have to be consulted. ALIUS Bulletin n7 (2024) aliusresearch.org/bulletin Changeux - Consciousness: from molecular to social scales 6 Now, there is a particular issue about experimentation with drugs since many of them are addictive - there is a loss of control taking place where use leads to abuse - and, even more, some of these drugs at a single dose may have irremediable effects on the brain (i.e., MPTP and Parkinson). Should it mean that the use of any of these drugs has to be forbidden? First of all, there is the use of drugs by shamans and native populations practicing shamanism. My view is that this traditional use of drugs is well controlled and, for millennia, was not shown to create irremediable alterations of consciousness. It might be accepted under well-defined social conditions. The dangers of uncontrolled use of drugs are well established (tobacco, alcohol, cocaine, heroin, etc.), and, by all means, the general population should be protected against it. The severity of the constraints imposed should be dictated by the dangerousness of the drug. This becomes particularly important when some of these dangerous drugs are planned to be used as medicaments against well-defined brain diseases. Classical, among many examples, is the use of cannabinoids on intractable pain, the use of nicotine against Gilles de la Tourette disease or Ritalin against ADHD despite its close relationships with cocaine. There is a dual role of drugs which should be examined seriously. A scientific debate about the clinical effect of the considered drug should take place to identify the positive effects together with the negative ones. An eventual prescription of any of such drugs should thus follow a prudence rule: the evaluation of the balance benefit-risks in the treatment of the disease. Moreover, each individual case has its own history and should be dealt with case by case. In light of your collaboration with artists and your work on the neural mechanisms of aesthetics, how do you believe altered states of consciousness, whether induced by psychoactive substances or other means, influence the cognitive and neural processes underlying the creation and perception of art? What implications might these altered states have for our understanding of aesthetic experience and creativity from a neuroscientific perspective? ALIUS Bulletin n7 (2024) aliusresearch.org/bulletin Changeux - Consciousness: from molecular to social scales 7 A number of important artists were taking drugs in their creative process. Famous examples are Henri Michaux with mescaline or Joan Mitchell with alcohol. There are also the Huichols from Mexico who created beautiful carpets inspired by their hallucinations during their pilgrimage to Wirikuta, during which they took peyotl to "communicate with the gods". "Psychedelic art" developed in the late 1960s as a counterculture, featuring highly distorted or surreal visuals and bright colors, conveying psychedelic experiences. I am not so much attracted to this kind of art, preferring the rational/emotional spirit of Poussin or Matisse. My view is that authentic artists do not need drugs to be creative, even if a few of them may receive some incentive from them. The works of art are artifacts, human productions, created for intersubjective communication. They use "symbolic forms" distinct from language, art is a mode of non-verbal communication of emotional states, imagination \& rational experience, under the constraints of "rules of art". Lev Vygotsky (1971) describes the art process as "emotional thinking". Upon contemplation, the work of art shows an esthetic efficacy: an "explosive" union of emotions \& imagination, referred to as catharsis by Vygotsky, which mobilizes conscious \& non-conscious processes. In extreme cases, there is the so-called Stendhal effect. I quote Stendhal: "I was already in a sort of ecstasy ... absorbed in the contemplation of sublime beauty, I saw it up close, I touched it so to speak ... Leaving Santa Croce, I had a heartbeat ... life was exhausted in myself, I walked with the fear of falling". Should we say that Stendhal describes some altered states of consciousness? In some cases, it may look like a pathological reaction of the brain, which needs medical help (a few cases per year at the Uffizi). Most of the time, one may view the esthetic catharsis as a physiological phenomenon, a particular case of ignition. Ignition is an experimentally well-documented process that manifests itself through the sudden activation and invasion of the global neuronal workspace and monitors conscious access (Mashour et al., 2020). Thus, the hypothesis is that aesthetic experience mobilizes a particular type of ignition, an "esthetic ignition," some kind of unforeseen global resonance reward within the GNW. ALIUS Bulletin n7 (2024) aliusresearch.org/bulletin Changeux - Consciousness: from molecular to social scales 8 The hypothesis is that aesthetic experience mobilizes a particular type of ignition, an "esthetic ignition," some kind of unforeseen global resonance reward within the GNW. The different components of the work of art (e.g., in a painting) are brought together globally in the conscious workspace. The phenomenon is supramodal. It concerns innate dispositions - as in the recognition of shape, color, space, and movement, but also of figures, their dispositions, meanings, and emotional valence - epigenetically acquired through a long process of selection of synapses during the lifetime experience of the participant as stored long-term memory of events and emotions, personal history, education that was received, and the experience of other artists' works) (Changeux et al., 1973). Esthetic ignition would be a naturally occurring physiological process - mobilizing jointly "reasons together with emotions" - that may occur in many cultures. Also, chills frequently accompanying a musical performance may be taken as a sign of esthetic ignition. Blood \& Zatorre (2001) have recorded parallel increase/decrease of activation in defined brain territories, many of them being concerned by drug addiction. Does esthetic experience deserve the qualification of an "altered" state of consciousness? Why not? This is to be further discussed! Generative neural networks have been increasingly used in artistic production. What is your opinion on the drawbacks and synergies that can operate between human and artificial intelligence when it comes to art and creativity? Some of the AI so called "artistic productions" might be of interest, and even sometimes give some "catharsis". We are under conditions which resemble that of ChatGTP: you recover what you have used to feed the system, yet sometimes with unexpected organizations but often these are systematic bric-a-brac compositions of cliches. On the other hand the real human artist attempts to convey to the spectator/auditor the expression of his inner feelings, what one may call his "personality", and transmit his own personal message. ALIUS Bulletin n7 (2024) aliusresearch.org/bulletin Changeux - Consciousness: from molecular to social scales 9 Art history enters into the actual style of the artist together with his own personal history, evolutive origins and the culture he/she belongs to. All that arises from the internalization in his brain of multiple experiences, education, unexpected encounters, original news... and of course the artistic context in which he/she is living. Nothing transcendental or irreducible in all that but the multiple intricate chances of epigenetic encounters, unpredictable chains of social circumstances, which make each individual human brain unique. There is no theoretical obstacle in principle with AI but machines simply miss an extraordinary intermingling between brain organization and cultural complexity, a theory of mind and individual history, among many other things which all are essential in artistic creation. To end on a positive note, of course AI can be fruitfully used as a tool for creation. These are the plethora of new techniques that the artist may practice to express him/herself. These include, in music, computer assisted composition, the development of new musical instruments, of new categories of sounds but also of methods for treating and manipulating sounds on line. In visual arts, there is the management of new forms, the treatment of large numbers of data. Machines simply miss an extraordinary intermingling between brain organization and cultural complexity, a theory of mind and individual history, among many other things, which all are essential in artistic creation. How do you perceive the role of artificial intelligence and machine learning in advancing our understanding of consciousness? AI and machine learning offer new and very efficient tools for scientific inquiry, from molecular biology to ecology, of course, and up to brain and consciousness. Apart from this general practical issue, the question is: to what extent do authentic homologies exist in how the artificial machines and the human brain are constructed to process information and eventually generate consciousness? ALIUS Bulletin n7 (2024) aliusresearch.org/bulletin Changeux - Consciousness: from molecular to social scales 10 There is a section of brain-inspired AI referred to as neuromorphic that attempts to develop silicon hardware resembling brain neurons and circuits. It has not yet developed and given results as outstanding as deep learning and ChatGTP with traditional AI. The principal aims of AI and machine learning are not to reproduce how the brain works but to simulate brain functions without necessarily any explicit reference to the brain. Spectacular successes have been achieved. However, there are intrinsic fundamental differences between machines and the brain that we cannot underestimate. Among them, AI machines benefit from a speed of processing orders of magnitudes faster (<speed of light) than the propagation of signals in the brain (<speed of sound). Also, they can embrace numbers of items with orders of magnitude larger than the human brain can do. Last, if the artificial simulation of conscious processing might have some successes in the future, silicons are not proteins and they shall miss (among many other features) the pharmacology of the real brain and the action of drugs... on consciousness! In light of your groundbreaking work on the neural foundations of individual consciousness, how do you conceptualize the social dimension of consciousness? Guillaume, you know that better than I do. I don't see a gap between individual and social consciousness since all of us know that the human brain is both rational and social. The social dimension has simply been missing in the past decades, and the time has come to reveal its importance. There is an urgent need for neuro-molecular conceptualization of social relationships such as theory of mind, decision-making, attachment, social bond, language, reasoning, science, art, ethics. The conceptualization should be directly inspired from the past work on individual consciousness yet with an original description of the relationships between individual consciousnesses. I have always been fascinated by the theories of Kropotkine (1902) on the importance of mutual aid in the competition between social groups, which according to him, are at the origins of Homo sapiens biological evolution. Yet the relationship between eventual genetic determinants of social life and biological evolution is still highly enigmatic. ALIUS Bulletin n7 (2024) aliusresearch.org/bulletin Changeux - Consciousness: from molecular to social scales 11 There is an urgent need for neuro-molecular conceptualization of social relationships such as theory of mind, decision-making, attachment, social bond, language, reasoning, science, art, ethics. Considering this social dimension, does the idea of collective consciousness extend to suggest a form of consciousness that emerges from social groups or societies, akin to concepts of collective intelligence? What would be the implications for understanding societal behaviors and decision-making processes? I have not been thinking about "collective consciousness" up to now. In any case, it should not be considered a metaphysical issue here. Since Durkheim (1912), it is usually considered as the set of shared beliefs, ideas, and moral attitudes uniting members within the society, creating solidarity, strengthening social bonds. There is, of course, an essential cultural component and, thus, a diversification of societies as a consequence of its epigenetic transmission. As a negative feature, it may lead to communitarianism and be the origin of multiple conflicts. When we see wars plaguing the world around us, we may realize that we are missing a collective consciousness that would share common universal values between conscious representatives of Homo sapiens. Would a stronger contribution of scientists improve the world situation? Not sure! But a science of collective consciousness should certainly develop! What do you think are the most promising areas for future research in altered states of consciousness? The most promising area is simply to further investigate, by similar methods, the neuro-molecular mechanisms of altered states of consciousness at the different levels or scales of brain physical organization (including the molecular level). This should include, in particular, the multiple levels of consciousness itself, such as minimal consciousness, recursive consciousness, self-consciousness, reflective consciousness, and theory of mind. These develop in the course of postnatal development in higher organisms, humans in particular, and in the course of biological evolution (see Lagercrantz \& Changeux 2009, 2010). ALIUS Bulletin n7 (2024) aliusresearch.org/bulletin Changeux - Consciousness: from molecular to social scales 12 I already mentioned the different modes of relationship with the environment i.e. body vs. outside world: physical, social, and cultural. Also, it is important to distinguish genetic from epigenetic conditions. All these are promising areas. The abundance of theories may be disorienting, and look useless, yet often they are complementary to each other and for all these reasons, there is the importance of open and unbiased scientific debate! You have often advocated for interdisciplinary approaches in neuroscience. How do you think fields like computational psychiatry and social neuroscience can contribute to our understanding of altered states of consciousness? How do you see the role of international collaborations, like ALIUS, in shaping the future of consciousness research? Althusser, in Philosophie et philosophie spontanee des savants (1967), criticized the slogan of "inter-disciplinarity" as leading to an internal fusion (or confusion) of disciplines that would then lose their identity. Personally, together with him, I prefer the term multi-disciplinarity. This means that representatives of different disciplines work together without losing the expertise they practice in their discipline, on the opposite. For instance, working on mental diseases like autism requires not only psychiatric expertise for the diagnosis but also human genetics, experimental psychology, brain imaging, developmental anatomy and chemistry, brain neurotransmitters and their receptors, synaptic epigenesis, at their most sophisticated levels. Nowadays, such expert multi-disciplinarity is needed to approach conscious processing successfully. I may take this opportunity to emphasize the importance of the too-often-forgotten social sciences such as anthropology or ethnology. In this respect, the late Martin Fortier reached a remarkable level of multi-disciplinarity, uniting basic sciences and social sciences. In his studies of the use of hallucinogenic drugs in the Shamanic practices of native populations of South America. He was able to combine expertise in molecular pharmacology, behavioral sciences, and anthropology, all that in his single scientist brain. Yet, multidisciplinarity is most often achieved through collaborations between several experts from different labs working cooperatively on the same topic even if conflicts between "disciplinary egos" sometimes happen. ALIUS Bulletin n7 (2024) aliusresearch.org/bulletin Changeux - Consciousness: from molecular to social scales 13 Collaborations are crucial to approach consciousness research, and international ones, such as ALIUS, are sometimes even more appropriate. What advice would you give to young scientists embarking on research in the complex field of consciousness? My first piece of advice to a young scientist is "to follow your own taste", as the poet Francis Ponge was saying, independently of fashion, financial pressure, scientific environment. As Andre Lwoff was saying: select a "good boss" who understands your intentions. Then, try to build a theoretical hypothesis - even simple-minded (better mathematical) - about your projects and develop the experimental tools to test it (see Volzhenin et al., 2022). In the field of consciousness, the efficient approaches are often multidisciplinary. Don't hesitate to learn about experimental techniques and theories in different labs - together with your own lab. This is a simple way to establish collaboration between labs. Avoid far-fetched theories and too complicated methods. Be simple minded and friendly. Good luck! My first piece of advice to a young scientist is "to follow your own taste" independently of fashion, financial pressure, scientific environment. References Althusser, L. (1967) Philosophie et philosophie spontanee des savants (1967) Maspero Paris Baars, B. (1988) A Cognitive Theory of Consciousness, NY: Cambridge University Press Blood, A.J., Zatorre, R.J. (2001) Intensely pleasurable responses to music correlate with activity in brain regions implicated in reward and emotion. PNAS. https://doi.org/10.1073/pnas.191355898 ALIUS Bulletin n7 (2024) aliusresearch.org/bulletin Changeux - Consciousness: from molecular to social scales 14 Changeux, J.P., Christopoulos, A. (2016) Allosteric modulation as a unifying mechanism for receptor function and regulation. Cell. https://doi.org/10.1016/j.cell.2016.08.015 Changeux J.P., Courrege P., Danchin, A. (1973) A theory of the epigenesis of neuronal networks by selective stabilization of synapses. PNAS. https://doi.org/10.1073/pnas.70.10.2974 Dehaene, S., Changeux J.-P. (1991) The Wisconsin Card Sorting Test: theoretical analysis and modeling in a neuronal network. Cereb Cortex. 1: 62-79. https://doi.org/10.1093/cercor/1.1.62 Dehaene, S., Kerszberg, M., Changeux, J.P. (1998) A neuronal model of a global workspace in effortful cognitive tasks. PNAS https://doi.org/10.1073/pnas.95.24.14529. Dehaene, S., Changeux, J.P. (2011) Experimental and theoretical approaches to conscious processing. Neuron. 70:200-27. https://doi.org/10.1016/j.neuron.2011.03.018. Dehaene, S., Sergent, C., Changeux, J.P. (2003). A neuronal network model linking subjective reports and objective physiological data during conscious perception. PNAS. 100 :8520-5. https://doi.org/10.1073/pnas.1332574100. Durkheim, E. (1912) : Les Formes elementaires de la vie religieuse : le systeme totemique en Australie, Paris, Felix Alcan Kropotkine, P. (1902) Mutual Aid: A Factor of Evolution London Lagercrantz, H., Changeux, J.P. (2009) The emergence of human consciousness: from fetal to ALIUS Bulletin n7 (2024) aliusresearch.org/bulletin Changeux - Consciousness: from molecular to social scales 15 neonatal life. Pediatr Res. 65:255-60. https://doi.org/10.1203/PDR.0b013e3181973b0d. Lagercrantz H, Changeux JP. (2010) Basic consciousness of the newborn. Semin. Perinatol. 34 :201-6. https://doi.org/10.1053/j.semperi.2010.02.004. Mashour, G.A., Roelfsema, P., Changeux, J.P., Dehaene, S. (2020) Conscious Processing and the Global Neuronal Workspace Hypothesis. Neuron. 105:776-798. https://doi.org/10.1016/j.neuron.2020.01.026. Volzhenin, K., Changeux, J.P., Dumas, G. (2022) Multilevel development of cognitive abilities in an artificial neural network. PNAS https://doi.org/10.1073/pnas.2201304119. Vygotsky L (1971) The Psychology of Art (English translation by MIT Press). ALIUS Bulletin n7 (2024) aliusresearch.org/bulletin\par
\endgroup

\ifdefined\ALIUSIssueBuild
\clearpage
\expandafter\endinput
\fi
\end{document}
